\documentclass[11pt,a4paper]{article}
\usepackage[margin=2.4cm]{geometry}
\usepackage{amsmath}
\usepackage{booktabs}
\usepackage{tabularx}
\usepackage[hidelinks]{hyperref}
\usepackage{parskip}
\usepackage[T1]{fontenc}
\sloppy

\title{\textbf{Satisfied, Not Solved}\\[2mm]
\large What a fairness fix actually does to real decisions}
\author{Dheirav}
\date{Findings summary --- \today}

\newcommand{\jargon}[1]{\emph{#1}}

\begin{document}
\maketitle

\begin{quote}
\small \textbf{Terminology.} Technical names are given in italics the first time they
appear, immediately after the plain description. A full glossary is in
\texttt{00-plain-english.pdf}.
\end{quote}

\section{The problem}

A program decides who gets a mortgage. It approves \textbf{79\%} of one group of applicants
and \textbf{57\%} of another. That gap is what people mean by algorithmic bias.

The standard fix is to instruct the program: \emph{approve both groups at the same rate}
(\jargon{demographic parity}). It works --- on the standard benchmark dataset the gap
between the two rates falls from $0.186$ to $0.018$, for less than two percentage points of
accuracy.

\textbf{But there are two completely different ways for a program to obey that
instruction:}

\begin{enumerate}
\item Approve \emph{more} people from the group being turned down.
\item Turn down \emph{more} people from the group getting through.
\end{enumerate}

Option 2 makes the numbers equal without helping anybody. It has a name in the ethics
literature: \jargon{levelling down}.

\textbf{The score that certifies the fix cannot tell these apart.} It measures only whether
the two rates match --- not how they came to match, and not whether the total number of
approvals went up or down.

\section{What was already known, and what was not}

\textbf{Known.} A 2023 paper by Mittelstadt, Wachter and Russell showed that these fixes
mostly do option 2, and argued it is the normal case rather than bad luck. It is in their
title: levelling down ``by default''.

\textbf{Also known, and I got this wrong at first.} That same paper \emph{also} proposes the
remedy --- their \S6 requires that every group receives at least some minimum approval rate.
I originally believed my version of this was new. It is not. I found that by reading their
paper properly, and corrected my write-up. What is different about mine is narrower: it
works without the program ever being told which group an applicant belongs to, and it is
tested on data the programs had never seen, across 26 datasets --- where they report one
dataset, scored on the same data they trained on.

\textbf{Also known: how to measure who paid.} A February 2023 paper by Ferry and
colleagues, \emph{When Mitigating Bias is Unfair}, sets out a framework for auditing what
one of these fixes does to individuals --- how many people it moves, in which direction,
and what it does to the overall approval rate. Those are the three things counted
throughout this work, so the measurement is theirs and I say so. What is mine is the result
of running it widely: they audit one model against one baseline, and never ask whether the
direction is a property of the data rather than of the method. It is.

\textbf{Also known, and published five months before I reached it.} A March 2026 paper,
\emph{Fairness May Backfire: When Leveling-Down Occurs}, proves that when the program cannot
see group membership --- the setting studied here --- the effect is distribution-dependent
and can go either way. I found this by running a novelty check after the fact, not before
starting. \textbf{The conditionality is therefore not my claim.} What that paper does not
have is any experiment: zero datasets appear in it.

\textbf{Also known: the silence of the score itself.} That an equalised score does not say
which of the two options produced it is reported in the literature as well --- a 2023 paper
on fairness across combinations of groups observes that levelling down ``often goes unnoticed
in the overall performance of the model'', and a 2023 auditing framework exists because the
aggregate score reports neither how many people were affected nor in which direction. Section
1 above is that known problem restated, not a claim of mine.

\textbf{Not known: \emph{when} it happens, in practice.} Their paper never discusses how generous the
system is overall. Their only remark near it is that the benchmark dataset is ``more than
75\% negatively labelled'' --- said in passing about accuracy cost, never connected to which
of the two options you get. They were working in exactly the regime where their conclusion
holds, and did not ask whether it would reverse elsewhere.

\textbf{The nearest miss.} A 2024 paper by Goethals, Delaney, Mittelstadt and Russell studies
how the cost of fairness changes with the available budget, and calls this ``a factor
overlooked in previous evaluations''. Two of those four authors also wrote the 2023 paper.
But in their setup the number of approvals is \emph{fixed in advance} as a budget, so the
total cannot move and the effect reported here is invisible by construction. The group best
placed to find this looked at the same dimension and did not make the connection.

\section{The main finding}

\textbf{Levelling down is not what these fixes do. It is what they do when the system is
stingy.}

Count what fraction of \emph{everyone} gets approved --- both groups together. Call it the
\textbf{approval rate} (\jargon{selection rate}).

\begin{itemize}
\item Approval rate below roughly \textbf{0.3}: the fix takes approvals away.
\item Approval rate above roughly \textbf{0.6}: the fix hands approvals out.
\item In between, it flips.
\end{itemize}

\subsection{How this was shown}

Two datasets always differ in dozens of ways, so comparing them proves nothing. Instead,
\textbf{one dataset was used and exactly one thing was changed.}

The dataset predicts whether someone earns more than \$50{,}000. That figure is an arbitrary
choice by whoever built the benchmark. Raise it and fewer people qualify; lower it and most
do. \textbf{Same people, same information about them, same groups --- only the finish line
moves.} Automated test code verifies that the people, the features and the groups really are
identical between runs.

\begin{center}
\small
\begin{tabular}{lrrr}
\toprule
Income cutoff & Approval rate & Change in total approvals & Lost per one gained \\
\midrule
\$100{,}000 & 0.030 & $-29.71\%$ & 22.03 \\
\$70{,}000  & 0.099 & $-22.05\%$ & 2.18 \\
\$50{,}000  & 0.252 & $-2.34\%$  & 1.14 \\
\$30{,}000  & 0.598 & $+0.98\%$  & 0.83 \\
\$20{,}000  & 0.760 & $+0.80\%$  & 0.75 \\
\$10{,}000  & 0.890 & $+0.08\%$  & 0.89 \\
\bottomrule
\end{tabular}
\end{center}

The last column is how many people lost an approval for each person who gained one. At the
top of the table, 22 people lose out for every one who benefits. At the bottom, more people
gain than lose.

\subsection{It holds up everywhere it was checked}

\begin{center}
\small
\begin{tabularx}{\linewidth}{lX}
\toprule
\textbf{Check} & \textbf{Result} \\
\midrule
A second population & Same pattern, more strongly \\
A different way of grouping people & Four populations, all agree \\
\textbf{Fourteen runs, on four new populations, made \emph{after} the prediction was written down} &
  All consistent; correlation $+0.901$ \\
Real mortgage decisions, where no arbitrary cutoff exists & Lands exactly where predicted \\
Two datasets it was never tuned on & Both placed on the correct side \\
\bottomrule
\end{tabularx}
\end{center}

\subsection{The obvious objection, answered in advance}

Someone will say: this is just the well-known fact that the two \emph{groups} have different
underlying rates.

It is not, and the difference matters. That objection is about a gap \emph{between two
groups}. This finding is about the \emph{overall level} for everybody together --- a
different quantity.

The two do move together, unavoidably: if almost nobody or almost everybody is approved,
there is no room for a gap between groups. So the analysis \textbf{holds the between-group
gap fixed and re-checks}, and the relationship gets \emph{stronger}, not weaker
($+0.801 \rightarrow +0.980$ in one population, $+0.964 \rightarrow +0.994$ in another).
That check is the core of the argument, not a footnote to it.

\section{How this sits against the theory}

Checking my measurements against that 2026 paper produced three results.

\textbf{Their conditions hold on none of my 26 datasets.} Their theorem is written over the
extrema of a ratio that diverges wherever its denominator approaches zero, and every real
dataset has such points. The two regions it compares always overlap, by large margins. The
conditions are sufficient rather than necessary, and as stated the theorem cannot be applied
to data --- which is consistent with the paper containing no experiments.

\textbf{Their direction is nonetheless right.} Relaxed from a strict separation to a
comparison of the two regions, it predicts correctly in 24 of 26; a median version gets 25.
Both misses are the two smallest effects in the set.

\textbf{And the approval rate is a proxy for their structural quantity}, at
\textbf{r = +0.935}. The two rules agree with each other on 25 of 26 datasets. Theirs
requires the joint distribution of the model's scores and group membership; mine requires
last year's approval rate.

So their mathematics says \emph{why}, this work says \emph{how to tell}, and the two agree
having been arrived at independently --- which is better evidence than either alone.

\section{Exactly what levelling down is}

There is a small piece of arithmetic that makes the whole thing precise, and it appears not
to be written down anywhere.

The fix puts both groups on one shared approval rate. Before the fix, the overall rate was
just the average of the two groups' rates, weighted by how many people are in each. So:

\begin{center}
\fbox{\parbox{0.9\linewidth}{\centering
The total number of approvals is unchanged \textbf{exactly when} the shared rate lands at
that size-weighted average.}}
\end{center}

Levelling down is nothing more mysterious than the program picking a shared rate
\emph{below} that point. This needs no assumptions and held to within $0.076$ percentage
points across all fourteen runs --- four fresh populations --- that were done only after
this was written down.

\section{What I could not work out}

\textbf{Why the flip happens.} I derived a candidate explanation from theory, wrote the
prediction and its pass mark down in advance, and tested it on fourteen runs across four
fresh populations, all done
afterwards. It passed the test I set --- and was then beaten by a far cruder rule
(``approval rate above one half? then it's the good kind''). So the derivation added nothing
and is reported as a failure.

The lesson is recorded with it: \textbf{I set a pass mark but never set a rival to beat}, so
passing did not mean what it looked like it meant.

\textbf{How big the effect is.} The approval rate tells you the direction, not the size. One
population loses $2.34\%$ where another loses $20.5\%$ at almost the same approval rate. A
candidate explanation for that was tested and refuted.

\section{Supporting findings}

\textbf{Combinations of groups get missed.} A fix applied to one characteristic leaves
people at the \emph{intersection} of two (say, Black women) worse off than any group the fix
was told about --- and worse than the benchmark dataset suggests ($9.0\times$ against
$13.2\times$). \textbf{This is not my observation.} Kearns et al.\ (2018) showed the failure
mode; Maheshwari et al.\ (2023) showed that the harm at the intersection is the larger one
and that it does not show up in a model's overall performance figures. What I add is the
scale --- ten populations --- and the condition that the effect needs a sizeable minority to
appear at all.

\textbf{The standard explainability tool does not reveal what happened.} A widely used method
for showing which information a program relied on reports a $+151\%$ shift after the fix.
Three explanations were proposed and none survived; most of the apparent shift turns out to
be an accounting artifact between two near-identical pieces of information. \textbf{The
transferable lesson: this tool can move enormously without the program's behaviour changing
much.}

\textbf{On small datasets the fix is drowned by chance.} Below roughly 2{,}500 test cases,
re-running with a different random split moves the numbers more than the fix itself does.

\textbf{The recommended alternative is worse.} \emph{Minimax group fairness} is proposed in
the literature as the answer to levelling down. On mortgage data it destroyed $10\%$ of all
approvals, at 46.7 lost per one gained --- on the very data where the ordinary fix
\emph{created} approvals. Its objective minimises the worst group's \emph{error rate}, and
the group with the most errors is not the group being treated worst.

\section{Why any of this matters}

Published figures for the domains where these fixes are deployed:

\begin{center}
\small
\begin{tabular}{lrl}
\toprule
Decision & Approval rate & Direction \\
\midrule
R\'esum\'e screening (applicant to interview, 2024) & 0.02--0.03 & takes opportunities away \\
Elite university admission & 0.04--0.10 & takes opportunities away \\
University admission, all four-year US institutions & 0.66--0.73 & hands them out \\
US mortgage lending (HMDA 2023 aggregate) & $\sim$0.84 & hands them out \\
\bottomrule
\end{tabular}
\end{center}

\textbf{They span the entire range and sit on opposite sides of the crossover.} My earlier
claim --- that these domains are nearly all stingy, so the harmful case is the default ---
was asserted and never checked, and it is wrong. The corrected version is stronger:

\begin{center}
\fbox{\parbox{0.92\linewidth}{\centering
\textbf{Two organisations can deploy the identical constraint, in good faith, and get
opposite effects on the number of opportunities granted --- with identical fairness
reports.}}}
\end{center}

That is why the direction cannot be assumed in either direction, and why a rule computable
in advance is worth having. It also validates the data used here: the HMDA populations sit
at 0.758 and 0.808 against a national aggregate of $\sim$0.84, so they are representative of
the domain rather than an unusually permissive slice.

\section{How carefully this was done}

Every result is the average of five runs with different random splits. The program is never
allowed to see the protected characteristic. All measures are written from their definitions
and checked against a standard library on every run.

Where a result could have gone either way, the prediction and its pass mark were written and
committed \emph{before} the experiment ran, and the order is provable from the version
history.

That discipline is why the following can be reported honestly: \textbf{three of my own
predictions failed}; two of my own proposed explanations were disproved by my own
experiments; and three claims were corrected after the fact --- including my claim that the
remedy was new, which the literature had got to first.

\section{The question I need answered}

A theory paper published in March 2026 proves this effect can go either way, and contains no
experiments. I have the experiments, a rule computable from a historical approval rate, and
the finding that their conditions are satisfied on none of my 26 datasets.

Is that an empirical paper worth publishing on its own --- or does it now read as a follow-up
to someone else's theorem?

\vfill
\noindent\rule{\textwidth}{0.4pt}\\
{\small Every number here is regenerated from stored results by an automated check that
fails if a document and its data disagree. The cited papers are in \texttt{papers/}; the
full research record is 16 documents in \texttt{research/docs/}.}

\end{document}
